\chapter{Complexity — Uniform Tableau}
%%% generated by the blueprint pipeline from TCSlib.Complexity.CircuitComplexity.UniformTableau
\section{Overview}

Declarations of the Lean module \texttt{TCSlib.Complexity.CircuitComplexity.UniformTableau}.

\section{Declarations}

\begin{theorem}[Polynomial-time construction of a delimited concatenation]
\lean{Complexity.polyTimeComputable_tabInput}
\label{Complexity.polyTimeComputable_tabInput}
\leanok
\uses{Complexity.PolyTimeComputable, Complexity.PolyTimeComputable.append, Complexity.polyTimeComputable_const}
\difficulty{3}
Let $m,T,z$ be polynomial-time computable functions from bit strings to bit strings.
Then the function sending each bit string $x$ to the concatenation of $m(x)$, the bit
$0$, $T(x)$, the bit $0$, and $z(x)$ is polynomial-time computable.
\end{theorem}

\begin{definition}[Time-bounded tabulation circuit family]
\lean{Complexity.tabFamily}
\label{Complexity.tabFamily}
\leanok
\uses{BoolCircuit.DAGCircuitFamily, Complexity.cfgTab}
Given a deterministic Turing machine $M$ and natural numbers $C$ and $d$, define the
circuit family whose circuit on inputs of length $n$ is the tabulation circuit for $M$
with empty initial tape contents and time bound $(C+1)(n+1)^d$.
\end{definition}

\begin{theorem}[Polynomial-time uniformity of the tableau family]
\lean{Complexity.tabFamily_isPUniform}
\label{Complexity.tabFamily_isPUniform}
\leanok
\uses{BoolCircuit.DAGCircuitFamily.IsPUniform, Complexity.PolyTimeComputable.comp, Complexity.UTab.polyTimeComputable_tabEmit, Complexity.UTab.tabEmit, Complexity.UTab.tabEmit_eq, Complexity.polyTimeComputable_const, Complexity.polyTimeComputable_id, Complexity.polyTimeComputable_polyUnary, Complexity.polyTimeComputable_tabInput, Complexity.tabFamily}
\difficulty{4}
For every finite-state Turing machine $M$ with Boolean input and output alphabets and
all natural numbers $C,d$, the circuit family whose $n$th circuit is the
configuration-tableau circuit for $M$ with time bound $(C+1)(n+1)^d$ is polynomial-time
uniform.
\end{theorem}

\begin{theorem}[Fan-in bound for the tableau circuit family]
\lean{Complexity.tabFamily_hasFaninTwo}
\label{Complexity.tabFamily_hasFaninTwo}
\leanok
\uses{BoolCircuit.DAGCircuitFamily.HasFaninTwo, Complexity.cfgTab_isFaninTwo, Complexity.tabFamily}
\difficulty{1}
Let $M$ be a finite-state Turing machine with Boolean tape symbols and
$C,d\in\mathbb{N}$. For each input length $n$, form the Boolean circuit encoding its
configuration tableau for $(C+1)(n+1)^d$ steps, with no fixed input prefix. Every gate
in every circuit of this family has at most two inputs; its input list has no
repetitions, and each NOT gate has exactly one input.
\end{theorem}

\begin{theorem}[Language of the tableau family]
\lean{Complexity.tabFamily_language}
\label{Complexity.tabFamily_language}
\leanok
\uses{BoolCircuit.DAGCircuitFamily.language, BoolCircuit.DAGCircuitFamily.mem_language_iff, Complexity.cfgTab_eval, Complexity.decide_emits_eq_of_computesInTime, Complexity.tabFamily, Turing.FinTM.ComputesInTime.mono, Turing.FinTM.DecidesInTime, Turing.MultiTapeTM.indicator}
\difficulty{4}
Let $M$ be a finite-state Turing machine with Boolean input and output alphabets, let
$L$ be a language of Boolean strings, and let $C,d$ be natural numbers. If $M$ decides
$L$ within time $C(n+1)^d$ on inputs of length $n$, then the language accepted by the
circuit family whose $n$th circuit is the configuration-tableau circuit for $M$ with
time bound $(C+1)(n+1)^d$ is exactly $L$.
\end{theorem}

\begin{theorem}[Polynomial-time-uniform circuits for polynomial-time languages]
\lean{Language.exists_isPUniform_of_mem_P}
\label{Language.exists_isPUniform_of_mem_P}
\leanok
\uses{BoolCircuit.DAGCircuitFamily, BoolCircuit.DAGCircuitFamily.HasFaninTwo, BoolCircuit.DAGCircuitFamily.IsPUniform, BoolCircuit.DAGCircuitFamily.language, Complexity.P, Complexity.mem_P_iff, Complexity.tabFamily, Complexity.tabFamily_hasFaninTwo, Complexity.tabFamily_isPUniform, Complexity.tabFamily_language}
\difficulty{4}
Every language of Boolean strings decidable in polynomial time has a family of Boolean
circuits with at most two inputs per gate, such that the circuit descriptions are
computable in polynomial time from the unary encoding of the input length and the family
accepts exactly that language.
\end{theorem}

\begin{theorem}[P-uniform circuit characterization of P]
\lean{Language.mem_P_iff_exists_isPUniform}
\label{Language.mem_P_iff_exists_isPUniform}
\leanok
\uses{BoolCircuit.DAGCircuitFamily, BoolCircuit.DAGCircuitFamily.HasFaninTwo, BoolCircuit.DAGCircuitFamily.IsPUniform, BoolCircuit.DAGCircuitFamily.language, Complexity.P, Language.exists_isPUniform_of_mem_P, Language.mem_P_of_isPUniform}
\difficulty{1}
A language $L$ over the Boolean alphabet belongs to polynomial time if and only if there
is a family of Boolean circuits that recognizes exactly $L$, has fan-in at most two, and
is polynomial-time uniform.
\end{theorem}
